\subsection{Integral class summation in a nonzero-Gram rational cell}
\label{sec:flag-generic-global-classes}

We retain the nonzero-Gram point, the two actual rational flags and
the local terms of Section~\ref{sec:flag-generic-local-amplitude}.
Write \(S\) for the intersection of their unipotent radicals and
\(H=U(J^\perp)\) for the pointwise stabilizer of their joint span \(J\).
The ambient signature is \((k+r,r+1)\); since \(J\) is nondegenerate
of dimension \(2r+2\) and contains an isotropic \(r\)-plane with a
perpendicular negative vector, it has signature \((r+1,r+1)\).
Thus \(H\) is compact at every infinite place. The notation \(S(F)\)
denotes rational isometries of the Hermitian space over \(E\).
The block calculation in the preceding subsection identifies \(S\)
over \(F\) with a product of restrictions of scalars of upper
unitriangular groups over \(E\).

At a finite place \(v\), use precisely the scaled frames of the
original \texttt{GFrames}: their \(i\)-th vectors are
\(c_v^{i-r}f_i,c_v^{i-r}f'_i\), where
\(c_vg=\delta_v\theta_v\), \(\delta_v\) generates the different of
\(F_v/\mathbb Q_p\), and \(\theta_v=1\) away from \(g\) and is a
uniformizer at the split level places. In particular
\(v(c_v)=v(\delta_v)\), since \(g\) is the product of the two
distinct split rational primes. The local spaces, inner lattices
and graphs are transported by this scaling, not changed.

\begin{lemma}[Integral local-to-global orbit product]\label{lem:generic-global-orbits}
Fix an occupied global datum \(d=(d_v)_v\), and one term lattice
\(\Lambda_0\) over it. Let \(\mathcal X(d)\) consist of the global
self-dual term lattices over \(d\), with the original local
admission conditions at every finite place. For every \(v\) let
\(\mathcal X_v(d_v)\) be the corresponding local term lattices.
Localization induces a bijection
\[
 S(F)\backslash\mathcal X(d)\ \longrightarrow\
       \prod_v' S(F_v)\backslash\mathcal X_v(d_v).
 \tag{GC1}
\]
The restricted product uses the orbit of \(\Lambda_{0,v}\).
If \(K_v^0\) is the original reference subgroup of \(S(F_v)\),
normalize its Haar measure to one. Then the original relative
global weight is
\[
 \mu(\Lambda)=
 \frac{[\Gamma^0:\Gamma^0\cap\Gamma_\Lambda]}
      {[\Gamma_\Lambda:\Gamma^0\cap\Gamma_\Lambda]}
      =\prod_v\mu_v(\Lambda_v),
 \qquad
 \mu_v(\Lambda_v)=
       \operatorname{vol}(S(F_v)_{\Lambda_v})^{-1},
 \tag{GC2}
\]
where \(\Gamma^0=S(F)\cap\prod_vK_v^0\) and
\(\Gamma_\Lambda=S(F)\cap\prod_v S(F_v)_{\Lambda_v}\).
\end{lemma}
\begin{proof}
A restricted family of local lattices patches uniquely to a global
lattice. Indeed choose a nonzero rational integer \(N\) bounding all
the finitely many changes between \(N\Lambda_0\) and
\(N^{-1}\Lambda_0\). The primary decomposition of
\(N^{-1}\Lambda_0/N\Lambda_0\) patches the prescribed local submodules
to exactly one submodule; duality, the lattice quotient conditions
and the level congruences can all be checked after localization.
This gives surjectivity of (GC1).

For injectivity, the local transporters between two lattices in the
same local orbits form nonempty open cosets of the local stabilizers.
They contain the identity at all but finitely many places. The
product is therefore a nonempty open subset of \(S(\mathbb A_{F,f})\).
In the block-unitriangular coordinates, ordinary additive
approximation in \(E\) gives density of \(S(F)\) in this finite
adelic group. Explicitly, to approximate finitely many entries
while remaining integral in the required fractional ideals
elsewhere, enlarge a common denominator supported at those
places and apply the Chinese remainder theorem to its fractional
ideal. The matrix entries and their inverse entries are polynomial
coordinates; the resulting approximation is an element of \(S(F)\).
It transports the lattices locally everywhere and hence globally.

For (GC2), put \(K=\prod_v S(F_v)_{\Lambda_v}\) and
\(K^0=\prod_vK_v^0\). The same density identifies the two rational
indices with the indices of \(K,K^0\) over their open intersection.
Their ratio is \(\operatorname{vol}(K^0)/\operatorname{vol}(K)\).
It is a finite product: the groups agree at almost every place,
where the two flags have integral corrected joint lattice of unit
determinant and the reference split coordinates apply.
\end{proof}

Only finitely many places, denoted \(P_*\), have \(H_v\ne1\),
where \(H_v\) is the local height of (GL2). At a place outside
\(P_*\), (GL4)--(GL7) give one orbit over each occupied datum, of
weight one. This orbit can be represented using the base frames:
the normalized frame class is the identity class, the joint
lattice is self-dual, and the complement splits off. Its phase
is one, using identity frame witnesses. At a place in \(P_*\)
the generic mass of a fixed datum is finite by (GL4)--(GL5).
The weights of the occupied frame classes and of each lattice
orbit are at least one, by the stabilizer inclusion in the proof
of (GL4). Thus there are finitely many generic orbits for each
fixed local datum.

For a generic lattice use the original finite phase
\(\omega_v\), the quotient of the two local simple-root
characters on witnesses. It is invariant under the common
unipotent. Consequently the phase-product formula and (GC1)--(GC2)
give the exact identities
\[
 \begin{split}
 A^0(d)&=\sum_{[\Lambda]\in S(F)\backslash\mathcal X(d)
                              \atop \Lambda\ {\rm generic}}
               \mu(\Lambda)\prod_v\omega_v(d_v,\Lambda_v)
                =\prod_{v\in P_*} A_v^0(d_v),\\
 m^0(d)&=\sum_{[\Lambda]\ {\rm generic}}\mu(\Lambda)
                =\prod_{v\in P_*}m_v^0(d_v).
 \end{split} \tag{GC3}
\]
A zero local generic mass makes both sides zero. These are
finite sums, not estimates obtained by bounding the number of
terms in a rational double coset.

\paragraph{Complement genera and the original row weights.}
For local occupied data \(d_v\), let \(H_{d_v}\) be their
stabilizers in \(H(F_v)\). Choose one global occupied datum
\(d^0\), normalize the Haar measures on \(H(F_v)\) by
\(\operatorname{vol}(H_{d_v^0})=1\), and write
\[
 V_0=\operatorname{vol}\bigl(H(F)\backslash H(\mathbb A_{F,f})\bigr),
 \qquad \mathrm{wt}_v(d_v)=1/\operatorname{vol}(H_{d_v}).
\]
Here the rational group is discrete in the finite adeles: a
compact set bounds denominators, while the real compactness
bounds the rational matrices at all infinite places, leaving a
finite set in the resulting arithmetic lattice. The mass is
finite by the usual finiteness of classes of definite Hermitian
lattices in a fixed genus, since the stabilizers are compact
open. This finiteness can also be applied to the complement
lattice of the reference datum and its finite residue graph;
passing to a finite-index subgroup changes the mass by a
finite factor. The actual reference genus mass is \(V_0\).

Consider a tuple of occupied local datum orbits, with the
reference orbit almost everywhere. It is realized globally:
choose term lattices in the prescribed local data and patch
them as in (GC1). All global data in the resulting genus are
parametrized by the double cosets
\[
          H(F)\backslash H(\mathbb A_{F,f})/\prod_v H_{d_v}.
\]
The same patching shows that every double coset gives occupied
data. Quotient integration on each double coset now gives
\[
 \sum_{[d']\in\mathrm{gen}(d)}
              \frac1{|H(F)_{d'}|}
           =V_0\prod_v\mathrm{wt}_v(d_v).
 \tag{GC4}
\]
No strong approximation for the definite unitary group is
asserted or required.

Here is the precise interpretation for the original row
classes. Their keys are the two fixed row lattices with their
relative quotient graph (the finite Levi label) and the
rational unipotent double coset. Keep the full set of original
row representatives. Fix a rational cell \(F,F'\), and denote
a term by an isometry \(\gamma:\Lambda_R\to\Lambda_L\) between
two such row lattices. An adapted mover \(\rho\) transports its
two flags to \(F,F'\) up to upper unitriangular changes.
The moved lattice is \(\rho\Lambda_L\).
Its two inner lattices are the images in the respective
isotropic quotients. Its graph is induced by
\[
       g\Lambda+(\Lambda\cap X)+(\Lambda\cap X'),
       \qquad X=\langle f_0,\ldots,f_{r-1}\rangle_E
 \tag{GC5}
\]
and the analogous \(X'\). In (GC5), a pair of quotient vectors
belongs to the graph when representatives have difference in
this module. Localizing this criterion and dividing the
isotropic graded generators by \(c_v^{r-i}\) gives exactly
the local \(\theta_v\)-graph; the factor \(g/\theta_v\) is a
local unit. Membership of the representatives and equality of
the graph can be checked locally, since all the modules in
(GC5) are lattices in their spans. Away from the level, the
difference condition is vacuous. At level places the
finite-adelic identification in the fixed-row congruence
calculation identifies equality of these graphs with equality
of the original \texttt{iotaOf} (not just equality of its
determinant). Thus the intrinsic datum in this cell is exactly
the two quotient lattices and the original row-label relation.

For a fixed key with representative \(\gamma_0,\rho\), write a
rational return as \(\gamma=n^{-1}\gamma_0m\), with \(n,m\) in
the left and right rational flag unipotents. Move its lattice
by \(\rho n\). The two inner lattices do not change, because
the unipotents act trivially on the respective isotropic
quotients. By (GC5), equality of the original row label is
precisely equality of the moved graph. Hence all moved
lattices of this key lie over its fixed datum \(d\).
Conversely, for any global term lattice over \(d\), its local
witnesses make it a local translate of \(\rho\Lambda_L\) by
the first flag unipotent, and of
\(\rho\gamma_0\Lambda_R\) by the second. Indeed the witness
splits the hyperbolic summands and leaves exactly the
prescribed inner lattice. The two rational unipotents are
split over \(F\) (one can use their triangular hyperbolic
coordinates and the trace-one correction as in
Section~\ref{sec:flag-reference-haar}); the same additive
approximation argument produces \(n,m\) globally with
\[
       \rho n\Lambda_L=\Lambda=\rho\gamma_0m\Lambda_R.
\]
The map \(n^{-1}\gamma_0m\) is thus integral between the original
rows. Its adapted frame conditions and level admission follow
locally from the witness conditions, so globally by lattice
intersection and the Chinese remainder theorem. It has the
given row label by (GC5). This constructs a term of the
original key. Two choices of \(n,m\) for one rational return
differ by a common unipotent after transport to the base cell.
Changing an integral return by either integral flag unipotent
does not change the moved orbit. Conversely if two moved
lattices differ by a common rational unipotent, its pullbacks
through the two movers preserve the corresponding original
row lattices, and hence are integral flag unipotents. This
proves the inverse implication. Therefore (GC1) counts
precisely the original integral terms of a fixed key. It
does not identify distinct terms inside the rational double
coset.

For a fixed global datum, let \(A_L,A_R\) be the finite
automorphism groups of its two definite inner lattices fixing
their last vectors, identified with the corresponding
\(\mathrm{autLz}\) of the original rows. The two row
identifications form torsors under these groups modulo
integral flag unipotents. To see that every identification
lifts, split off the isotropic ideal modules and their duals
in a self-dual lattice as in the fixed-row construction:
the middle lattice carries the given inner isometry and the
map on hyperbolic summands is the identity in the adapted
coordinates. Different choices of such a splitting act
trivially on the graded flag spaces and on the isotropic
quotient, hence differ by an integral flag unipotent.
The same argument gives an adapted rational realization
with the prescribed fractional flag ideals.

Two pairs of identifications give the same row/rational-cell
key if and only if their induced data are related by an
isometry of the two base flags. In fact equality of rational
double cosets provides the two rational unipotent
translations which, combined with the identifications, give
one simultaneous adapted isometry of both flags. Such an
isometry lies in \(S(F)H(F)\): on the two frame spans its
action is exactly a compatible pair of upper unitriangular
matrices, and its remaining action is an isometry of
\(J^\perp\). The common unipotent acts trivially on the
data, and equality of the row labels is equality of the
graph by (GC5). Hence the remaining isometry is precisely an
element of \(H(F)_d\). Conversely such an element preserves
both inner lattices and the graph, so its simultaneous change
of the two identifications preserves the key. This action is
free on the pair of identifications, because an element of
\(H\) acting trivially on an isotropic quotient is the
identity (the complement embeds in that quotient).
The number of keys over the datum class is therefore
\[
          |A_L|\,|A_R|/|H(F)_d|.
\]
Each key has squared row weight \(1/(|A_L||A_R|)\). Their
sum is exactly \(1/|H(F)_d|\). In particular, on any subset
of keys closed under the prescribed local datum relation, or
on any finite subcollection of it, the row-weight sum over
one tuple of local datum classes is bounded by the right side
of (GC4). This is the required genus-weight unfolding.
The argument is about keys, not about a single representative
for each key's integral terms.

\paragraph{The absolute arithmetic height and generic mass.}
Write \(f=[F:\mathbb Q]\), \(\rho=\operatorname{rd}(E)\)
(in this paragraph \(\rho\) is a number, not a flag mover),
and let \(D_l\) be the rational nonzero endpoint minors, with
\(G\) the rational joint Gram invariant. Set
\[
 N_G=\prod_{\tau:F\hookrightarrow\mathbb R}
                    |\tau(G)|\,\rho^{-2r(r+1)},\qquad
 N_D=\prod_{l\ {\rm endpoint}}\prod_\tau
             |\tau(D_l)|^2\,\rho^{-2l(2r+1-l)} .
\]
At a complex embedding the absolute square means the
embedding of \(D_l\overline{D_l}\); at every finite place
the same norm expression applies. Scaling the two frames
multiplies the \(l\)-th endpoint minor by
\(c_v^{-l(2r+1-l)}\). Both terms in the Gram invariant have
factor \(c_v^{-2r(r+1)}\). Since \(E/F\) is unramified at
finite places, \(N_F\mathfrak d_F=\rho^f\).
The ideal product formula thus gives
\[
 \prod_v q_v^{a_v}=g^{-2(r+1)f}N_G,\qquad
 \prod_v q_v^{d_v}=N_D,\qquad
 H_{\rm tot}:=\prod_vH_v=g^{-(2r+1)f}N_GN_D.
 \tag{GC6}
\]
These are the original \(H_p\) and \(HPt\) with all level
powers retained. In particular \(H_{\rm tot}\ge1\), and
\(N_D\ge1\) on an occupied point.

For each \(v\) choose the actual occupied reference pairing
class \(C_v\) used to define the generic mass; outside
\(P_*\) it is the unique class. Their frame witnesses can
be chosen as the localizations of one pair of rational
adapted frames \(u,u'\) up to linked integral changes.
Indeed the required local cosets in the two unitriangular
groups are open, and additive approximation in their
entries meets the finitely many specified cosets with
integrality relative to the frame ideals elsewhere.
Let \(M\) be the full lattice in the nondegenerate joint
space \(J\) with these corrected local generators; it exists
by the same patching argument. Define the literal gluing mass
\[
 {\cal G}(M)=
 \sum_{\Lambda=\Lambda^\#,\ \Lambda\cap J=M\ \bmod H(F)}
                       |H(F)_\Lambda|^{-1}.
\]
This is a finite sum, since the positive complement has
fixed discriminant module and there are finitely many
definite lattices of that discriminant up to isometry, with
only finitely many gluing maps. It is nonzero by the primitive
local density and patching.

At a fixed global datum \(d\) the mass of its generic terms
in the selected frame classes equals
\(\prod_v w(C_v)\) times the number of global lattices over
\(d\) having these prescribed frames and saturation \(M\).
This follows from (GC1)--(GC2) and the fixed-frame count
in (GL6): the lattices in the prescribed frames patch
bijectively with the local lattices in those frames.
Unfolding the complement orbits by (GC4) yields the exact
identity and its bound
\[
 ngen_{\rm tot}:=V_0\prod_{v\in P_*}ngen_v(C_v)
       =\left(\prod_{v\in P_*}w(C_v)\right){\cal G}(M)
       \le H_{\rm tot}^{s(s-1)}{\cal G}(M),\qquad s=r+1.
 \tag{GC7}
\]
Here \({\cal G}(M)\) is the source's \texttt{glueSum} for
the adapted rational pair, with its original fixed-frame
pointwise stabilizer weights. The product does not contain
a factor for the number of arbitrary adapted representatives.
This sharpens the intermediate source loss \(3s^2\).
Equation (GC7) is a reduction to a definite gluing mass,
not yet a comparison of that mass with the source's
normalized mean constant \(AA\).

\paragraph{Real pinned scalar and the phase.}
At a relevant nonzero-Gram point, write \(b(w)\) for the
number of strict upper entries within its consecutive
blocks (the source's \(\#\mathrm{FreeE}(w)\)).
The reference rational common unipotent has block entries
\(x_{ij}\in\mathfrak d^{-(j-i)}\), since the corresponding
scaled local entries are \(c_v^{j-i}x_{ij}\). The other
frame's action is obtained by the reversed inverse
adjoint on each block, and this preserves the same
integrality conditions. In the pinned real coordinates
the free entries are obtained by conjugating the second
frame's matrix by the upper triangular frame change and
scaling each height \(h=j-i\) by \(\rho^h\).
For clarity, if \(B=1+X\) is this block matrix, and \(U\)
is any upper unitriangular frame change, then
\((U^{-1}XU)_{ij}=X_{ij}\) plus a linear combination of
entries \(X_{ab}\) with \(b-a<j-i\).
This follows directly by expanding the two products, and
remains true after restriction to the original free
entries. The diagonal change from the prescribed real
flag basis to the unit coordinates gives the factor
\(\rho^{j-i}\); the other changes in the real frame have
diagonal one (unit complex factors have no effect on
the absolute Jacobian). This is the pinned chart, rather
than a new Haar measure on \(S(F_\infty)\).
Consequently the coordinate change has triangular Jacobian
one after those indicated diagonal scalings.

The lattice of a single free coordinate of height \(h\)
has complex embedding volume
\[
   2^{-f}\sqrt{|\operatorname{disc}E|}\,
   N_E(\mathfrak d^{-h})\,\rho^{2hf}
      =(\rho/2)^f.
\]
The original pinned Haar density for each free complex entry
at each real place is \(ka/(g\pi)=2/(\rho g)\).
Multiplication over the \(f\) real places and the \(b(w)\)
free entries gives \(g^{-fb(w)}\).
Successive reduction in order of entry height gives a
fundamental domain, since multiplication and conjugation
have triangular translation laws. Thus the product volume
also holds for the nonabelian group, including the empty
set of entries. Finally (GC2) or the common finite-index
formula gives the actual covolume
\[
         \operatorname{covol}(\Gamma_\Lambda)
                        =g^{-fb(w)}\mu(\Lambda).
 \tag{GC8}
\]
This calculation uses the reference root normalization of
the source's \texttt{IsPinnedR}. It does not equate an
arbitrary, not yet pinned quotient measure with that
normalization.

For the phase, put
\(s(A)=\sum_i A_{i,i+1}\) for an upper unitriangular
frame matrix, so \(s(AB)=s(A)+s(B)\).
Let \(U,U'\) be the two adapted changes of a moved term,
and \(\xi\) its actual integral admission witness. The
two compatible local frame matrices are \(U\) and
\(U'\xi^{-1}\) in the rational frames.
Local scaling multiplies their simple-root sums by
\(c_v\). The original local character is
\(\psi_{0,v}(x)=\psi_{E,v}(-x/\delta_v)\), evaluated
at the simple-root sum divided by \(\theta_v\).
The exact identity \(c_vg=\delta_v\theta_v\) therefore
gives
\[
 \omega_v=
 \psi_{E,v}\bigl((-s(U)+s(U')-s(\xi))/g\bigr).
\]
The global additive product formula, with the original
real sign convention, gives
\[
 \prod_v\omega_v
   =\psi_\infty(s(U))^{-1}\psi_\infty(s(U'))
                          \psi_\infty(s(\xi))^{-1},
 \qquad
 \psi_\infty(x)=\exp(2\pi i\,\operatorname{Tr}_{E/\mathbb Q}(x)/g).
 \tag{GC9}
\]
Equivalently, this follows by applying the product formula
to each of the three root sums; the finite characters
are trivial away from finitely many places. This keeps the
integral admission phase rather than making all terms
positive.

For a fixed row key and rational double coset with point
\(w\), the real orbital integral is the same operator
after transport by the two rational flag unipotents,
up to their simple-root characters and the covolume
of their common stabilizer. This follows by changing
variables in the two unipotent quotients: relevance is
exactly what makes their shared character quotient
trivial on \(S\), and integration over the arithmetic
stabilizer first gives its covolume. In the pinned
real quotient normalization the resulting operator is
\(J_w(\nu)\), with the original global scalar
\(\mathrm{scl}\,\mathrm{vol}(K_2)\). Unipotent translation
of the witnesses contributes the ratio of their two
characters. The factor from the original level character
is \(\psi_\infty(s(\xi))^{-1}\). By (GC9) their product
is the finite phase times a fixed unit depending only on
the base mover. More explicitly, if the original
base mover has matrices \(U_0,U'_0\), replacing them
by \(U_0 A,U'_0 A'\) multiplies the right side of (GC9)
by \(\psi_\infty(s(A))^{-1}\psi_\infty(s(A'))\), with
the new admission witness retained. This is precisely
the factor from the change of variables and its level
character. Absorb the fixed unit in \(T_w\), so
\(\|T_w\|_{\rm HS}=\|J_w(\nu)\|_{\rm HS}\).
Thus on this pinned cell each amplitude is
\(g^{-fb(w)}\mu(\Lambda)\prod_v\omega_v\,T_w\).
Summing the original weighted vectors over the generic
integral returns of this key gives, by (GC3),
\[
 \left\|\sum_{t\ {\rm generic\ in\ key}}\operatorname{vec}(t)
            \right\|^2
 =\mathrm{wt}(\mathrm{key})^2\,g^{-2fb(w)}
                  \|J_w(\nu)\|_{\rm HS}^2
                   \prod_{v\in P_*}|A_v^0(d_v)|^2.
 \tag{GC10}
\]
The Hilbert--Schmidt normalization is the original one.
This is an equality of the phase-sensitive sum, rather
than a Cauchy--Schwarz loss in the number of integral terms.
If the local generic set is empty both sides vanish.
If the real operator vanishes, (GC10) still holds as a
homogeneous equality, with no division by its norm.

Equations (GC3)--(GC10) provide the class-summation
and genus-weight components of the generic local input.
They apply to the original pinned real cell measures;
an additional comparison is needed wherever the selected
measure is only a bounded scalar multiple of these
measures. To finish \texttt{StepLocGenPinnedG}, one must
make that scalar comparison uniformly on the whole
family and identify \({\cal G}(M)\) with the explicit
mean mass, including the source's \(AA\) normalization.
The point-sum estimate in \texttt{StepThmBPinnedG}
is not a consequence of the arithmetic identities alone.
